\section{Sports Desk: Query Predict}
\label{sec:query-predict}

Query Predict (\texttt{qp}) is Emi's approach to \textbf{hand-written SQL
with generated, type-safe Go bindings}. You write plain \texttt{.sql}
files; Emi parses them, predicts the shape of each result row and the
parameters it needs, and generates a typed row struct and a
prepared-statement builder. There is no schema introspection, no
auto-migration, and no SQL synthesis --- the SQL you write is the SQL
that runs. This keeps the full power of SQL in your hands while still
giving compile-time safety over the columns and parameters a query
touches.

\begin{sidebar}{Three commands}
\textbf{emi qp:dir} --- scans a directory of \texttt{.sql} files, one
generated Go file per query, named from the SQL file.

\textbf{emi qp:gen} --- reads a Query Predict YAML document bundling
several named queries and a Go package name.

\textbf{emi qp:sql} --- compiles one query into executable Go and
transforms its metadata into pure SQL.
\end{sidebar}

\begin{lstlisting}
emi qp:dir --path ./queries-source --output ./output
\end{lstlisting}
For a query like \texttt{select name, id from users;}, Emi emits a typed
row struct, the original SQL as a constant, and a context/prepare helper:
\begin{lstlisting}
const BasicSelectSQL = `select name, id from users;`
type BasicSelectRow struct { Name string; Id string }
type BasicSelectContext struct {
  Filter, Having, Restriction string
  Params map[string]interface{}
}
\end{lstlisting}

\subsection{SQL helper functions}
A small set of marker functions can appear inside the SQL itself, letting
dynamic behaviour stay in the query while the generated code remains
predictable and typed:
\begin{sidebar}{Query Predict SQL helpers}
\scriptsize
\begin{tabular}{@{}p{0.42\linewidth}p{0.5\linewidth}@{}}
\texttt{field(expr, 'type', 'goName'?)} & annotate an expression's Go type (and optionally its field name) when it can't be inferred \\
\texttt{useval('name')} & a named value substituted from \texttt{ctx.Params} at runtime, safely escaped \\
\texttt{filter()} & replaced by the \texttt{Filter} clause from the context (default \texttt{1}) \\
\texttt{restriction()} & replaced by the \texttt{Restriction} clause (default \texttt{1}) \\
\end{tabular}
\end{sidebar}
\begin{lstlisting}
SELECT u.user_id,
  field(u.user_name, 'string') as UserName,
  field(COUNT(o.order_id), 'int64') AS total_orders
FROM users u LEFT JOIN orders o ON u.user_id = o.user_id
WHERE filter()
GROUP BY u.user_id, u.user_name
LIMIT useval('limit')
\end{lstlisting}
Here \texttt{field(..., 'int64')} tells the generator the aggregate
column is an \texttt{int64}, \texttt{filter()} is swapped for whatever
filter the caller supplies at runtime, and \texttt{useval('limit')} pulls
the \texttt{limit} entry out of \texttt{ctx.Params} --- Emi never invents
schema, migrations, or queries of its own; it only makes the boundary
between hand-written SQL and Go type-safe.
